% ======================================================================
\section{Problems and Standing Notation}\label{sec:problems}

Notation is that of the companion memo~\cite{memo1} (cited as M1): $H\in\mathbb{C}^{m\times n}$ is HSS on a
tree $\mathcal T$; a leaf $\tau$ has $n_\tau$ rows, $p_\tau$ columns, generators
$D_\tau,U_\tau,V_\tau$, and the non-leaf generators are $R,W,B$. Write
$\mathcal M(H,b):=H^+b$ for the output of the recursive minimum-norm solver, which M1
(Thm.~\ref{M1-thm:general}) proves exact for every full-row-rank HSS matrix with proper ranks
(M1, Assumption~\ref{M1-as:ranks}). In this memo $L$ denotes the weight, not the tree depth. We treat
\begin{align}
\text{(W)}\quad & x_W=\operatorname*{arg\,min}_{Hx=b}\|Lx\|_2, && H \text{ wide, full row rank},\label{eq:W}\\
\text{(T)}\quad & x_\lambda=\operatorname*{arg\,min}_{x}\ \|S(Hx-b)\|_2^2+\lambda^2\|Lx\|_2^2, && \lambda>0,\ H\text{ of any shape},\label{eq:T}
\end{align}
with weights that are \emph{diagonal or block diagonal conforming to the leaves}:
\[
L=\operatorname{diag}(L_\tau)_{\tau\ \text{leaf}},\quad L_\tau\in\mathbb{C}^{p_\tau\times p_\tau}\ \text{invertible};\qquad
S=\operatorname{diag}(S_\tau),\quad S_\tau\in\mathbb{C}^{n_\tau\times n_\tau}.
\]
The diagonal case $L=\operatorname{diag}(w)$, $w_j\neq0$, is the one used by IRLS and by
prior-weighted reconstructions. Note $\|Lx\|^2=x^*Mx$ with $M=L^*L$ SPD, so (W) is the familiar
$x_W=M^{-1}H^*(HM^{-1}H^*)^{-1}b$.

Two facts from M1 are used repeatedly.
\begin{itemize}[leftmargin=1.4em,itemsep=1pt]
\item[(F1)] (M1, Lemma~\ref{M1-lem:blockrow}) A leaf's block column is $H(:,J_\tau)=\big[\,F_\tau V_\tau^*\ \text{(other rows)};\ D_\tau\ \text{(own rows)}\big]$,
and its block row is $\big[\,D_\tau,\ U_\tau G_\tau\,\big]$. Operations on a leaf's columns act on
$(D_\tau,V_\tau)$ only, operations on its rows act on $(D_\tau,U_\tau)$ only.
\item[(F2)] (M1, Lemmas~\ref{M1-lem:orthogonal}--\ref{M1-lem:rank}, Table~\ref{M1-tab:chain}) Each operation of the solver either is a unitary
change of variables (feasible sets map onto each other, the norm is preserved) or removes variables
that enter the norm additively and are free (size reduction) or fixed by the constraints (forced
solve).
\end{itemize}

% ======================================================================
\section{Weighted Minimum Norm}\label{sec:weighted}

\begin{lemma}[change of variables]\label{lem:Wsub}
With $y=Lx$: $x_W=L^{-1}y^\star$, where $y^\star=(HL^{-1})^+b$ is the minimum-norm solution of
$(HL^{-1})y=b$.
\end{lemma}
\begin{proof}
\emph{Map:} $x\mapsto y=Lx$ is a bijection of $\mathbb{C}^n$. \emph{Feasible sets:}
$Hx=b\iff (HL^{-1})y=b$. \emph{Objective:} $\|Lx\|=\|y\|$. So (W) is the minimum-norm problem for
$HL^{-1}$ in the variable $y$.
\end{proof}

\begin{lemma}[$HL^{-1}$ is HSS on the same tree]\label{lem:Wstruct}
$H_L:=HL^{-1}$ has the same tree, ranks, and generators $U,R,W,B$ as $H$, with leaf generators
\[
D_\tau\ \to\ D_\tau L_\tau^{-1},\qquad V_\tau\ \to\ L_\tau^{-*}V_\tau .
\]
$H_L$ has full row rank, and it satisfies M1's Assumptions~\ref{M1-as:slack} and~\ref{M1-as:ranks} exactly when $H$ does.
\end{lemma}
\begin{proof}
$L^{-1}$ is block diagonal on the leaf column blocks, so it multiplies each block column
$H(:,J_\tau)$ on the right by $L_\tau^{-1}$. By (F1) that block column is $[F_\tau V_\tau^*;D_\tau]$, and
$F_\tau V_\tau^*L_\tau^{-1}=F_\tau(L_\tau^{-*}V_\tau)^*$. Nothing else is touched. Right multiplication by an
invertible matrix preserves the rank, and $n_\tau,p_\tau,\ell_\tau$ are unchanged.
\end{proof}

\begin{theorem}[weighted solver]\label{thm:W}
For (W) with leaf-conforming $L$, the vector $x=L^{-1}\mathcal M(HL^{-1},b)$ is $x_W$. In the
original variable, every operation of the solver either preserves the \emph{weighted} norm
$\|L\,\cdot\,\|$ or removes directions that only add to it, as in Table~\ref{tab:Wchain}.
\end{theorem}
\begin{proof}
By Lemma~\ref{lem:Wstruct}, $H_L$ is a full-row-rank HSS matrix, so M1 Thm.~\ref{M1-thm:general} gives
$\mathcal M(H_L,b)=y^\star$, and Lemma~\ref{lem:Wsub} gives $x_W=L^{-1}y^\star$.

For the step-by-step statement, run M1's chain (P0)--(P3) on $H_L$ in the variable $y$ and pull
each change of variables back to $x=L^{-1}y$. Since $L^{-1}$ acts leaf by leaf:
\begin{itemize}[leftmargin=1.4em,itemsep=2pt]
\item \textbf{Size reduction} on $H_L$: $[V_\tau^*L_\tau^{-1};D_\tau L_\tau^{-1}]\Omega_\tau=[\tilde V_\tau^*,0;\tilde D_\tau,0]$.
In $x$: $x_\tau=L_\tau^{-1}\Omega_\tau[\tilde y_\tau;w_\tau]$, so
$\|Lx\|^2=\sum_\tau\|\tilde y_\tau\|^2+\|w_\tau\|^2$. The discarded $w_\tau$ do not affect $Hx$, so the
optimum sets them to zero (M1 Lemma~\ref{M1-lem:orthogonal}). These are exactly the directions in which moving $x$
costs weighted norm and buys nothing.
\item \textbf{Row compression} $P_\tau^*$ is unchanged: it acts on rows (on $U_\tau$, $D_\tau$, $b$),
and $L$ acts on columns, so the two commute. It does not change the feasible set.
\item \textbf{Column decoupling} on $H_L$: $\tilde y_\tau=Q_\tau\hat y_\tau$. In $x$:
$x_\tau=L_\tau^{-1}\Omega_\tau[Q_\tau\hat y_\tau;0]$, and $\|Lx\|=\|\hat y\|$. This map is an isometry
from $(\hat y,\|\cdot\|)$ to $(x,\|L\,\cdot\|)$.
\item \textbf{Forced solve}: $X^{(1)}=\mathbf L^{-1}\hat b_1$ is the same for every feasible point, so it
contributes the constant $\|X^{(1)}\|^2$ to $\|Lx\|^2$ (M1 Lemma~\ref{M1-lem:forced}). Nonsingularity of the
triangular blocks follows from full row rank of $H_L$ (M1 Lemma~\ref{M1-lem:rank}).
\item \textbf{Recursion}: the reduced matrix $\bar H_L$ is HSS with one level fewer, has full row rank
and proper ranks (M1 Lemma~\ref{M1-lem:rank}, Sec.~\ref{M1-subsec:merge}). Induct.
\end{itemize}
Composing the maps gives $x=L^{-1}y^\star=x_W$.
\end{proof}

\begin{table}[h]
\centering\small
\begin{tabular}{@{}lll@{}}
\toprule
operation (on $H_L=HL^{-1}$) & map back to $x$ & invariant\\
\midrule
weight change & $x=L^{-1}y$ & $\|Lx\|=\|y\|$, \ $Hx=b\iff H_Ly=b$\\
size reduction & $x_\tau=L_\tau^{-1}\Omega_\tau[\tilde y_\tau;w_\tau]$ & $\|Lx\|^2=\|\tilde y\|^2+\|w\|^2\Rightarrow w=0$\\
row compression & rows $\times P_\tau^*$ & feasible set unchanged\\
column decoupling & $x_\tau=L_\tau^{-1}\Omega_\tau[Q_\tau\hat y_\tau;0]$ & $\|Lx\|=\|\hat y\|$\\
forced solve & $\hat y=[X^{(1)};X^{(2)}]$, $X^{(1)}$ fixed & $\|Lx\|^2=c+\|X^{(2)}\|^2$\\
recursion / reconstruction & as in M1 & M1 Thm.~\ref{M1-thm:general}\\
\bottomrule
\end{tabular}
\caption{Weighted version of M1's Table~\ref{M1-tab:chain}: each step preserves $\|L\,\cdot\|$ or removes
directions that only add to it.}
\label{tab:Wchain}
\end{table}

\begin{remark}[cost and implementation]
Forming $H_L$ touches only the leaf generators: $O(\sum_\tau p_\tau(n_\tau+\ell_\tau))=O(n\,r)$ work
for diagonal $L$, and $O(\sum_\tau p_\tau^3)$ for dense leaf blocks. The solve costs the same as for
$H$. In MATLAB: \texttt{minnorm(H,b,'Weight',w)} or \texttt{minnorm(H,b,'Weight',\{L\_1,\dots,L\_q\})},
which forms $H_L$ by editing the leaf generators ($D_\tau L_\tau^{-1}$ and $L_\tau^{-*}V_\tau$, by solves) and calls
M1's solver unchanged. In Python: \texttt{hssmn.weighted\_minnorm}.
\end{remark}

\begin{remark}[weights that are \emph{not} leaf-conforming]\label{rem:nonconform}
If $L$ couples columns of different leaves, then $HL^{-1}$ is no longer obtained by editing leaf
generators. Examples: a circulant or Toeplitz smoothing operator, or a weight that is diagonal
in a \emph{different} basis. Two practical cases from the application list are of this type:
spectral weights in seismic minimum weighted norm interpolation (MWNI~\cite{mwni}) and image-domain weights
in non-Cartesian MRI (Sec.~\ref{sec:apps}). When $L$ is itself HSS, $HL^{-1}$ is HSS of somewhat
higher rank, but forming it needs HSS inversion and HSS$\times$HSS products. The repo does not
have those yet, and this memo does not treat that case.
\end{remark}

% ======================================================================
\section{Tikhonov Regularization}\label{sec:tikhonov}

\subsection{Standard form via an augmented minimum-norm problem}

\begin{lemma}[augmentation]\label{lem:aug}
For $\lambda>0$ and any $H\in\mathbb{C}^{m\times n}$, let $[x^\star;s^\star]$ be the minimum-norm solution of
\begin{equation}\label{eq:aug}
\begin{bmatrix}H & \lambda I_m\end{bmatrix}\begin{bmatrix}x\\ s\end{bmatrix}=b .
\end{equation}
Then $x^\star=x_\lambda:=\operatorname*{arg\,min}_x\|Hx-b\|^2+\lambda^2\|x\|^2$, and $s^\star=(b-Hx_\lambda)/\lambda$.
\end{lemma}
\begin{proof}
\emph{Map:} for every $x\in\mathbb{C}^n$ there is exactly one $s$ with $(x,s)$ feasible, namely
$s=(b-Hx)/\lambda$. So $x\mapsto(x,(b-Hx)/\lambda)$ is a bijection from $\mathbb{C}^n$ onto the feasible set.
\emph{Objective:} on the feasible set,
$\|[x;s]\|^2=\|x\|^2+\lambda^{-2}\|b-Hx\|^2=\lambda^{-2}\big(\|Hx-b\|^2+\lambda^2\|x\|^2\big)$.
So the two minimizations are the same problem. Both have unique minimizers: the Tikhonov
functional is strictly convex, and \eqref{eq:aug} has full row rank.
\end{proof}
The closed form agrees: the minimum-norm solution of \eqref{eq:aug} is
$[H^*;\lambda I](HH^*+\lambda^2I)^{-1}b$, so $x_\lambda=H^*(HH^*+\lambda^2I)^{-1}b=(H^*H+\lambda^2I)^{-1}H^*b$.
This augmented formulation is classical (it is the ``damped'' problem of LSQR~\cite{lsqr}; see also~\cite{hansen}). The new point is
that it \emph{stays HSS on the same tree}, so M1's solver applies unchanged:

\begin{lemma}[the augmented matrix is HSS on the same tree]\label{lem:augstruct}
Order the unknowns leaf by leaf as $(x_{J_\tau},s_{I_\tau})_\tau$, i.e.\ apply the column permutation
$\Pi$ that puts the $n_\tau$ identity columns belonging to leaf $\tau$'s rows right after $\tau$'s
own columns. Then $H_a:=[H,\lambda I]\Pi$ is HSS on $\mathcal T$ with the same $U,R,W,B$ and
\[
D_\tau^{a}=\begin{bmatrix}D_\tau & \lambda I_{n_\tau}\end{bmatrix}\in\mathbb{C}^{n_\tau\times(p_\tau+n_\tau)},\qquad
V_\tau^{a}=\begin{bmatrix}V_\tau\\ 0_{n_\tau\times\ell_\tau}\end{bmatrix}.
\]
Moreover: (i) $H_a$ has full row rank for every $\lambda>0$ and every shape of $H$, including square
and tall; (ii) Assumption~\ref{M1-as:slack} of M1 holds at every leaf as soon as $\ell_\tau\le p_\tau$
($\ell_\tau+n_\tau\le p_\tau+n_\tau$), so even the legacy MATLAB file, which merges levels when the slack
is short, never collapses the leaf level of $H_a$ (the current solver never merges); (iii) $\Pi$ is a permutation, so it maps minimum-norm solutions to minimum-norm solutions.
\end{lemma}
\begin{proof}
The diagonal block of leaf $\tau$ in $H_a$ is $[H(I_\tau,J_\tau),\ \lambda I(I_\tau,I_\tau)]=[D_\tau,\lambda I]$.
An off-diagonal block between siblings $\alpha,\beta$ has columns $J_\beta$ of $H$ together with the identity
columns of $I_\beta$. Since $I_\alpha\cap I_\beta=\emptyset$, the identity contributes $0$, so the block is
$[\mathbf U_\alpha B_{\alpha\beta}\mathbf V_\beta^*,\ 0]=\mathbf U_\alpha B_{\alpha\beta}[\mathbf V_\beta^*,0]$. With
the leafwise interleaving, the zero columns sit under each leaf, which is $V_\tau^a=[V_\tau;0]$ with the
same $W$ above. (i): the columns of $\lambda I$ alone have rank $m$. (ii): direct. (iii): a permutation is
unitary (M1 Lemma~\ref{M1-lem:twosided} with $P=I$).
\end{proof}

\begin{theorem}[Tikhonov solver]\label{thm:T}
For $\lambda>0$ and any HSS matrix $H$ with proper ranks, $x_\lambda=\big[\Pi\,\mathcal M(H_a,b)\big]_{x\text{-part}}$.
In terms of the Tikhonov functional $J_\lambda(x)=\|Hx-b\|^2+\lambda^2\|x\|^2$, every operation of
the solver applied to $H_a$ preserves $J_\lambda$-optimality, as in Table~\ref{tab:Tchain}.
\end{theorem}
\begin{proof}
By Lemma~\ref{lem:augstruct}, $H_a$ is HSS with full row rank, so M1 Thm.~\ref{M1-thm:general} gives
$\mathcal M(H_a,b)=H_a^+b$. Undoing $\Pi$ gives the minimum-norm solution of \eqref{eq:aug}, and
Lemma~\ref{lem:aug} identifies its $x$-part with $x_\lambda$.

Step by step, with $z=(x,s)$ and $J_\lambda(x)=\lambda^2\|z\|^2$ on the feasible set (Lemma~\ref{lem:aug}):
\begin{itemize}[leftmargin=1.4em,itemsep=2pt]
\item \textbf{Size reduction} compresses the $(\ell_\tau+n_\tau)\times(p_\tau+n_\tau)$ stack
$\left[\begin{smallmatrix}V_\tau^*&0\\D_\tau&\lambda I\end{smallmatrix}\right]$. It mixes $x_\tau$ and $s_\tau$
(data-misfit and solution variables of the same leaf) by a unitary $\Omega_\tau$, and always
discards $p_\tau-\ell_\tau$ columns, \emph{even when the leaf of $H$ is square or tall}. Along the
discarded directions $z$ changes but $H_az$ does not, so $J_\lambda$ would only grow: set them to zero.
\item \textbf{Row compression} acts on $U_\tau$, $D_\tau^a$ and $b_\tau$. It is the same $P_\tau$ as for $H$
(the identity block becomes $\lambda P_\tau^*$, still part of $D_\tau^a$). It does not change the feasible set.
\item \textbf{Column decoupling} is a unitary change of $z_\tau$, so it preserves $\|z\|$, hence $J_\lambda$.
\item \textbf{Forced solve}: the forced unknowns are the same for every feasible $z$. They add a
constant to $\|z\|^2$, i.e.\ to $J_\lambda/\lambda^2$. Nonsingularity: M1 Lemma~\ref{M1-lem:rank}, since $H_a$ has full
row rank.
\item \textbf{Recursion}: the reduced matrix has full row rank and one level fewer. Induct.
\end{itemize}
Each step is either a unitary change of $z$ or removes components of $z$ that are free or fixed and
enter $\|z\|^2$ additively. Composed, the steps map the minimizer of the last problem to the minimizer
of $J_\lambda$.
\end{proof}

\begin{table}[h]
\centering\small
\begin{tabular}{@{}lll@{}}
\toprule
operation & variable / map & invariant\\
\midrule
augmentation & $x\mapsto z=(x,(b-Hx)/\lambda)$ & $\lambda^2\|z\|^2=J_\lambda(x)$ on $\{H_az=b\}$\\
interleave $\Pi$ & permutation of $z$ & $\|z\|$\\
size reduction & $z_\tau=\Omega_\tau[\tilde z_\tau;w_\tau]$ & $\|z\|^2=\|\tilde z\|^2+\|w\|^2\Rightarrow w=0$\\
row compression & rows $\times P_\tau^*$ & feasible set\\
column decoupling & $\tilde z_\tau=Q_\tau\hat z_\tau$ & $\|z\|$\\
forced solve & $\hat z=[Z^{(1)};Z^{(2)}]$, $Z^{(1)}$ fixed & $\|z\|^2=c+\|Z^{(2)}\|^2$\\
recursion & as in M1 & M1 Thm.~\ref{M1-thm:general}\\
\bottomrule
\end{tabular}
\caption{The chain behind Theorem~\ref{thm:T}.}
\label{tab:Tchain}
\end{table}

\subsection{General form, data weights, limits}
\begin{proposition}[general form]\label{prop:general}
For (T) with leaf-conforming $S$ and invertible leaf-conforming $L$, put $\tilde H=SHL^{-1}$,
$\tilde b=Sb$. Then $x_\lambda=L^{-1}\tilde y_\lambda$, where $\tilde y_\lambda$ is the standard-form
Tikhonov solution for $(\tilde H,\tilde b)$. $\tilde H$ is HSS on the same tree with
$D_\tau\to S_\tau D_\tau L_\tau^{-1}$, $U_\tau\to S_\tau U_\tau$, $V_\tau\to L_\tau^{-*}V_\tau$, so
Theorem~\ref{thm:T} applies.
\end{proposition}
\begin{proof}
With $y=Lx$ the objective is $\|\tilde H y-\tilde b\|^2+\lambda^2\|y\|^2$, a bijective change of
variables. The structure follows from (F1) on rows and on columns, as in Lemma~\ref{lem:Wstruct}.
\end{proof}
$S$ may be singular: a 0/1 mask that drops corrupted data rows keeps $[\tilde H,\lambda I]$ at full row rank.
\emph{Limits.} As $\lambda\to0$, $x_\lambda\to H^+b$: the minimum-norm solution for wide full-row-rank $H$,
the least-squares solution for tall full-column-rank $H$. So the Tikhonov path connects to M1
continuously. As $\lambda\to\infty$, $x_\lambda\approx\lambda^{-2}H^*b$.
\emph{Conditioning.} $\kappa(H_a)^2=(\sigma_1^2+\lambda^2)/(\sigma_{\min}^2+\lambda^2)$ for wide or square $H$,
and $(\sigma_1^2+\lambda^2)/\lambda^2$ for tall $H$. For wide and square $H$, regularization improves the
conditioning of the problem the solver actually sees. For tall $H$, $\kappa(H_a)\to\infty$ as
$\lambda\to0$ although $x_\lambda$ tends to the least-squares solution, whose conditioning is that of
$H$; there $\kappa(H_a)u$ is a loose bound for the error in $x$ (Section~\ref{subsec:Tnum}).
\emph{Cost.} Augmented leaves have $p_\tau+n_\tau$ columns, but the size reduction brings them back to
$\ell_\tau+n_\tau$, the width a leaf of $H$ has after its own size reduction when its slack suffices. Only that first QR is
larger, by at most the factor $1+n_\tau/p_\tau$; every later operation has the same sizes as for $H$.
So a Tikhonov factorization costs little more than a minimum-norm one (measured: $1.0$--$1.25\times$).
Each $\lambda$ needs a new factorization, since $D_\tau^a$ depends on $\lambda$. A $\lambda$ sweep of
length $K$ costs $K$ factorizations, each linear in $n$.
In MATLAB: \texttt{[x,r] = tikhonov(H,b,lambda,'Weight',L,'DataWeight',S)}, with $L$ and $S$ vectors or
cells of leaf blocks. In Python: \texttt{hssmn.tikhonov}.
